\section{Stability of a maintained interface}\label{sec:stability}

\subsection{A stationary force balance is only the starting point}
Let $(\bm y_\star,\bm a_\star)$ be a candidate maintained state and command.
For an autonomous slow-state or acoustic-envelope description, define
$\bm\chi=\bm y-\bm y_\star$ in local coordinates satisfying the conservation
constraints. Denote command perturbation by $\delta\bm a=\bm a-\bm a_\star$,
and the dynamic state and command derivative operators by $\mathsf A$ and
$\mathsf B$, respectively. Linearization gives
\begin{equation}
 \dot{\bm\chi}=\mathsf A\bm\chi+\mathsf B\delta\bm a.
 \label{eq:linear-state}
\end{equation}
The operators $\mathsf A,\mathsf B$ are derivatives of the dynamic equations,
not derivatives of a surface-fitting routine. The state includes the necessary
fluid, thermal, actuator and acoustic memory variables. A steady phasor field
represents a periodic physical wave; if an envelope reduction is unjustified,
stability must instead be studied about that periodic state.

For a finite-dimensional autonomous reduction, negative real parts of all
eigenvalues imply local linear asymptotic stability. For the continuum problem,
point spectrum alone is insufficient: suitable evolution-operator bounds are
needed. Define $\mathcal T(t)$ as the linear perturbation propagation operator,
$M_s\geq1$ as a dimensionless amplification bound, and $\alpha_s>0$ as a decay
rate. A stronger useful condition is
\begin{equation}
 \norm{\mathcal T(t)}\leq M_s e^{-\alpha_s t},\qquad t\geq0.
 \label{eq:semigroup-bound}
\end{equation}
The norm uses physically stated state scales. A large $M_s$ permits substantial
transient growth despite eventual decay. Acceptability additionally
depends on the observation $\mathsf C_{\rm obs}\bm\chi$, where
$\mathsf C_{\rm obs}$ is the linearized application-specific observation operator, and on
the amplitude of admitted disturbances. Numerical convergence, shape accuracy
and dynamical stability are three different claims.

\subsection{Linearize the actual base flow in three dimensions}
If the maintained surface carries a steady streaming flow $\bm U_\star$,
then $\bm U_\star\cdot\nn=0$ on the surface does not imply
$\bm U_\star=0$ in the fluids. For constant-density mean momentum, the
convective derivative contributes
\begin{equation}
 \delta[(\bm U\cdot\nabla)\bm U]
   =(\bm U_\star\cdot\nabla)\delta\bm U
     +(\delta\bm U\cdot\nabla)\bm U_\star.
 \label{eq:streaming-linearization}
\end{equation}
Both terms, incompressibility, perturbed viscous stress, moving-interface
conditions and the derivatives of acoustic forcing belong in $\mathsf A$.
Temperature-dependent coefficients and source loading contribute further
terms when present. A spectrum obtained after setting this measured or
computed base circulation to zero is the spectrum of a different problem.

The cylindrical container admits angular perturbations of every integer order
$m$. If the complete base state and boundary laws are rotationally invariant,
$\mathsf A$ commutes with rotations and can be decomposed into sectors
proportional to $e^{\mathrm i m\theta}$. Stability requires control of all
admissible sectors, including $m\ne0$, not just the axisymmetric sector.
An asymmetric base state couples these orders and needs the corresponding
three-dimensional linearized operator. Conservation constraints, contact-line
motion and legitimate symmetry modes must be handled before interpreting
zero eigenvalues. A finite angular truncation also needs a bound or a
convergence argument for its omitted modes.

\subsection{Capillary stiffness and acoustic feedback}
In the quiescent, constant-$\sigma$ reduction, define the shape operator
$\bm B_\Gamma=\Gs\nn$ and its squared norm $|\bm B_\Gamma|^2$, equal to the sum
of squared principal curvatures. Let $\Delta_\Gamma=\Gs\cdot\Gs$ denote the
surface Laplacian. Under normal displacement $\eta$,
\begin{equation}
 \delta\kappa=-(\Delta_\Gamma+|\bm B_\Gamma|^2)\eta,
 \qquad \delta\Phi=(\nn\cdot\nabla\Phi)\eta.
 \label{eq:curvature-variation}
\end{equation}
At fixed actual commanded variable, denote the complete shape derivative of
acoustic normal load by $D_\Gamma\Pi[\eta]$. The restoring-pressure operator,
before projection onto volume-preserving variations, is
\begin{equation}
 \boxed{\mathcal K_{\rm eff}\eta
 =-\sigma(\Delta_\Gamma+|\bm B_\Gamma|^2)\eta
    +\Delta\rho(\nn\cdot\nabla\Phi)\eta
    -D_\Gamma\Pi[\eta].}
 \label{eq:effective-stiffness}
\end{equation}
The pressure-multiplier variation is removed on that constrained space. Boundary
terms and contact-line perturbations must be included for the chosen wetting
model. Thermal variations add derivatives of $\sigma$, density and other
properties. This formula shows the relevant question: when the interface moves,
does the resulting acoustic load counteract or reinforce the displacement?
Matching the load at one geometry supplies no answer to that derivative.

If fluid inertia and acoustic memory are negligible, let $\mathcal M_\Gamma$
be the viscous mobility mapping outward traction to interface normal velocity,
with its constraints and boundary conditions. Then the creeping-flow limit is
\begin{equation}
 \dot\eta=-\mathcal M_\Gamma\mathcal K_{\rm eff}\eta.
 \label{eq:stokes-stability}
\end{equation}
For passive Stokes flow the constrained mobility is positive in the reciprocal
work pairing. A coercive self-adjoint restoring operator can then give an
energy-based stability criterion. Driven, dissipative acoustic feedback need
not be self-adjoint or derive from a scalar shape potential; positivity of a
guessed ``capillary plus acoustic energy'' is not a general stability proof.

With memory retained, let $s\in\C$ be a Laplace frequency,
$\mathcal Z_h(s)$ the hydrodynamic restoring/dynamic impedance, and
$\mathcal K_a(s)$ the linearized acoustic load response to surface displacement.
The homogeneous perturbation condition is
\begin{equation}
 [\mathcal Z_h(s)-\mathcal K_a(s)]\eta=0.
 \label{eq:memory-stability}
\end{equation}
Both operators include their relevant boundary and state coupling. In a
particular modal approximation, hydrodynamics may be represented by
$\mathcal Z_h(s)=s^2\mathsf M+s\mathsf D+\mathsf K_h$, where
$\mathsf M,\mathsf D,\mathsf K_h$ are the modal mass, damping and passive
stiffness matrices. Frequency-independent damping is an approximation, not a
universal viscous free-surface law; capillary-wave dispersion depends on the
viscous regime \cite{denner2016}. Eliminating acoustic state instantaneously
replaces $\mathcal K_a(s)$ by its zero-frequency value and can miss an unstable
oscillatory mode.

For a periodically maintained state of period $T_p$, let $\mathcal M_p$ denote
the linear evolution over one period. In a finite-dimensional periodic
reduction, its Floquet multipliers must lie strictly inside the unit circle
after genuine symmetry modes are treated appropriately. A continuum statement
also requires operator control. Pulse trains, modulation and beat forcing may
therefore need a periodic analysis even when their mean load is correct.

\subsection{Feedback must have the necessary authority and observations}
Let the measured perturbation be
$\bm z_{\rm obs}=\mathsf C\bm\chi+\bm\nu$, where $\mathsf C$ is the sensor
operator and $\bm\nu$ measurement noise. The existence of a desired equilibrium
does not imply that its unstable modes can be controlled or observed. Local
stabilization requires, at minimum, authority over the unstable modes and
sufficient information to estimate those that need feedback. Delays, saturation
and finite sensing bandwidth belong to the control model.

One precise necessary test uses a left unstable mode. Suppose an isolated
mode $\psi$ of the adjoint state operator satisfies
$\mathsf A^*\psi=\bar\lambda\psi$, with $\Rea\lambda\geq0$ and
$\mathsf B^*\psi=0$. In the usual complex pairing, its amplitude
$a_\psi=\inner{\psi}{\bm\chi}$ satisfies
\begin{equation}
 \dot a_\psi=\lambda a_\psi
       +\inner{\mathsf B^*\psi}{\delta\bm a}
       =\lambda a_\psi.
 \label{eq:uncontrollable-mode}
\end{equation}
No linear feedback through the admitted source operator can make this mode
decay exponentially. For an infinite-dimensional boundary-control problem,
the adjoint pairing requires its appropriate trace and admissibility domain;
the statement is for modes on which that pairing is defined. Similarly, an
unstable state mode invisible to the measurements obstructs its asymptotic
estimation by a detectable linear observer. These tests connect source support
and sensing to stability, independently of a target-fitting algorithm.

For illustration only, consider real, scaled finite-dimensional coordinates
and full-state feedback $\delta\bm a=-\mathsf K_c\bm\chi$. Here
$\mathsf K_c$ is a feedback gain, distinct from the traction matrices
$\bm K_i$. Let $\mathsf P_L\succ0$ be a symmetric Lyapunov matrix. A
sufficient linear decay condition is
\begin{equation}
 (\mathsf A-\mathsf B\mathsf K_c)^T\mathsf P_L
 +\mathsf P_L(\mathsf A-\mathsf B\mathsf K_c)
 \preceq-2\alpha_s\mathsf P_L.
 \label{eq:lyapunov}
\end{equation}
With $\mathsf X=\mathsf P_L^{-1}$ and $\mathsf Y=\mathsf K_c\mathsf X$, this
becomes the linear matrix inequality
\begin{equation}
 \mathsf A\mathsf X+\mathsf X\mathsf A^T
 -\mathsf B\mathsf Y-\mathsf Y^T\mathsf B^T
       +2\alpha_s\mathsf X\preceq0,\qquad\mathsf X\succ0.
 \label{eq:feedback-lmi}
\end{equation}
This established construction \cite{boyd1994} is convex for fixed
$\mathsf A,\mathsf B,\alpha_s$ and the stated full-state model. Joint design
of geometry, nonlinear acoustics and output feedback does not inherit that
convexity. A controller satisfying it still needs a compatible estimator,
actuator headroom and a nonlinear region of attraction containing the intended
formation endpoint and disturbances.

\subsection{What is needed to extend a linear certificate locally}
A conditional estimate makes the limitation of a linear result concrete.
In scaled state coordinates suppose the closed-loop nonlinear perturbation
equation has the form
\begin{equation}
 \dot{\bm\chi}=\mathsf A_c\bm\chi+\mathcal N(\bm\chi),
 \qquad \norm{\mathcal N(\bm\chi)}\leq c_N\norm{\bm\chi}^2
              \quad\text{for }\norm{\bm\chi}\leq r_0.
 \label{eq:nonlinear-remainder}
\end{equation}
Here $\mathsf A_c=\mathsf A-\mathsf B\mathsf K_c$, $c_N\geq0$ is a
remainder bound, and $r_0>0$ is a neighborhood radius. Assume local
well-posedness there and a bounded positive Lyapunov operator $\mathsf P_L$
for which \cref{eq:lyapunov} holds as a quadratic-form inequality on the
generator domain, replacing matrix transpose by the Hilbert adjoint when
using an operator state space. Let $p_-,p_+>0$ satisfy
$p_-\norm{\bm\chi}^2\leq V_L\leq p_+\norm{\bm\chi}^2$, where
$V_L=\inner{\bm\chi}{\mathsf P_L\bm\chi}$. For classical solutions,
\begin{equation}
 \dot V_L\leq-2\alpha_s V_L
                   +2p_+c_N\norm{\bm\chi}^3.
 \label{eq:nonlinear-lyapunov-bound}
\end{equation}
If $c_N>0$, set
$r_*=\min\{r_0,\alpha_s p_-/(2p_+c_N)\}$; if $c_N=0$, take $r_*=r_0$.
For $\norm{\bm\chi}\leq r_*$ the positive term is at most
$\alpha_s V_L$, hence $\dot V_L\leq-\alpha_s V_L$. Initial data with
$V_L(0)<p_-r_*^2$ remain in that neighborhood and decay, as long as the
stated solution and control constraints remain valid.

This is a local conditional estimate, not a claim that the free-boundary
equations satisfy the assumed quadratic remainder in an arbitrary norm.
In a continuum problem, regularity and possible derivative loss require
care in choosing that state space. Formation must end inside a verified
admissible neighborhood, and source saturation or disturbances require
their own bounds. A negative finite-matrix spectrum by itself supplies none
of these nonlinear hypotheses.
